\newcommand{\NominalTable}{%
\begin{table}[htbp]
\centering
\caption{Nominal inputs for the reported calculations.}\label{tab:parameters}
\begin{tabular}{ll}\toprule
Quantity & Value\\\midrule
Substrate radius $R$ & $12.7$ mm\\
Substrate thickness $L$ & $5.0$ mm\\
Beam radius $w_0$ & $3.81$ mm ($0.3R$)\\
Full-plane absorbed power $P_{\rm abs}$ & $0.110$ W\\
Thermal conductivity $k$ & $1.38$ W\,m$^{-1}$\,K$^{-1}$\\
Density $\rho$ & $2203$ kg\,m$^{-3}$\\
Specific heat $c_p$ & $772$ J\,kg$^{-1}$\,K$^{-1}$\\
Ambient temperature $T_\infty$ & $300$ K\\
Convection coefficient $h_{\rm conv}$ & $10$ W\,m$^{-2}$\,K$^{-1}$\\
Emissivity $\epsilon$ & $0.85$\\
Linearized total loss $h_{\rm eff}$ & $15.2054$ W\,m$^{-2}$\,K$^{-1}$\\
Thermo-optic coefficient $\gamma_T$ & $9.6\times10^{-6}$ K$^{-1}$\\
Reference wavelength; refractive indices & $1064$ nm; $n_i=1$, $n_s=1.44963$\\
Entrance-pupil radius $a$ & $0.9R$\\\bottomrule
\end{tabular}
\par\smallskip\begin{minipage}{.96\linewidth}\footnotesize $\gamma_T$ is a retained nominal value tabulated at $643.8$ nm, not a wavelength-specific measurement at $1064$ nm. The intercepted disk power is obtained by integrating the prescribed flux on the finite face.\end{minipage}
\end{table}}

\newcommand{\SteadyFieldTable}{%
\begin{table}[htbp]
\centering
\caption{Steady temperature rises and sampled OPD-PV.}\label{tab:steadyfields}
\begin{tabular}{lrrrr}\toprule
$\theta$ & Front center (K) & Midplane center (K) & Rear center (K) & OPD-PV (nm)\\\midrule
$0^\circ$ & 11.921 & 7.656 & 6.716 & 166.503\\
$30^\circ$ & 11.336 & 7.537 & 6.656 & 168.771\\
$60^\circ$ & 9.303 & 6.959 & 6.314 & 154.105\\\bottomrule
\end{tabular}
\par\smallskip\begin{minipage}{.96\linewidth}\footnotesize OPD-PV uses the saved $201\times201$ physical transmitted-pupil grid. Continuous numerical extrema are reported separately in Supplement 1; FEM comparisons use their own discrete common domain.\end{minipage}
\end{table}}

\newcommand{\TransientRefinementTable}{%
\begin{table}[htbp]
\centering
\caption{Largest transient FEM refinement changes across temperature and OPD metrics (percent).}\label{tab:refinement}
\begin{tabular}{lrrrr}\toprule
$\theta$ & Transverse & Axial & Time & Joint\\\midrule
$0^\circ$ & 0.104616 & 0.029993 & 0.025763 & 0.125259\\
$30^\circ$ & 0.099671 & 0.029973 & 0.025837 & 0.118577\\
$60^\circ$ & 0.088451 & 0.029925 & 0.026017 & 0.105834\\
Target & $0.5$ & $0.5$ & $0.1$ & $0.5$\\\bottomrule
\end{tabular}
\par\smallskip\begin{minipage}{.96\linewidth}\footnotesize Maxima include the positive observation times from $0.0025$ to $75$ s. The $0.001$ s frame is recorded separately as a diagnostic; the exact zero initial condition and all-step physical checks remain independent requirements.\end{minipage}
\end{table}}

\newcommand{\ValidationTable}{%
\begin{table}[htbp]
\centering
\caption{Optical differences between FEM and the semi-analytical calculation (percent).}\label{tab:comparison}
\begin{tabular}{lrrr}\toprule
$\theta$ & Steady $E_V$ & Transient $E_V$ & Transient $E_{W,\infty}$\\\midrule
$0^\circ$ & 0.004220 & 0.071263 & 0.090254\\
$30^\circ$ & 0.000523 & 0.054394 & 0.085118\\
$60^\circ$ & 0.005440 & 0.046821 & 0.082570\\\bottomrule
\end{tabular}
\par\smallskip\begin{minipage}{.96\linewidth}\footnotesize Steady values use the finest FEM mesh. Transient values are maxima over sampled times from $2.5$ ms to $75$ s and use the FEM candidate examined by the separate refinements. All denominators are the semi-analytical PV on the common pupil. Supplement 1 retains full temperature and optical statistics, including $1$ ms diagnostics.\end{minipage}
\end{table}}
